% !TeX root = ../../Book2.tex
% 纲要标题：### 20.6* 非交换 Gröbner 基与 Diamond 引理
% 纲要位置：计划纲要.md，第 1731 行。

% 纲要要点（供目录核对）：
%
% * 自由结合代数中的可容许字序、重写规则及包含和重叠歧义；下降次序须与两侧连接相容，两个最小失败例的实际商在正文命题完整证明
% * Diamond 引理与正规字；分别证明终止、首步选择无关及正规映射核等于关系理想，不把交换 Buchberger 判据原样照搬
% * PBW 型代数及可解型多项式代数中的有效约化；区别单次约化终止、全部歧义可解、规则有效匹配与补全终止，有限展示不自动给出有限 Gröbner 基
% * 建议先读第一卷附录 F.1—F.3 的交换 Gröbner 基；第三卷第8章发展计算应用、附录 F 处理局部标准基；正文用同一理想的两种字序证明有限基依赖字序，原题改为乘法、中心与单性计算，不把维数数据当作乘法结构
% ---
%

\OptionalSection{非交换 Gröbner 基与 Diamond 引理}{b2:sec:2006}

第一卷附录F.1—F.3讨论交换多项式的除法、首项理想和Buchberger 算法\footnote{布赫贝格尔（Bruno Buchberger，1942-），奥地利数学家，提出 Gröbner 基算法。}。本节在这些概念的背景下重新定义非交换字的次序与重写规则，说明字母顺序带来的差别。

\ProseParagraphBreak
正规式不仅要能够算出，还须证明不同计算次序给出相同结果。本节在自由结合代数的字上建立一个首一重写版本的 Diamond 引理。规则下降保证每次约化过程终止，歧义可解才保证不同选择得到相同正规式。要把它用于算法，还须说明规则如何匹配、系数如何计算；有限条原始关系也不保证补全后仍只有有限条规则。这些是不同的要求。


\subsection{可容许字序、重写规则与歧义}
\label{b2:subsec:2006:01}

令 \(X\) 为有限字母集，\(X^*\) 为全部有限字，空字记为一。自由结合代数 \(k\langle X\rangle\) 以这些字为基，乘法为连接。

\begin{definition}\label{b2:def:admissible-word-order}
设 \(X\) 为有限字母集，\(X^*\) 为由 \(X\) 构成的全部有限字的集合，空字记为 \(1\)。若 \(X^*\) 上的全序 \(<\) 为良序，以空字为最小元，且对任意字 \(a,b,u,v\in X^*\)，\(u<v\) 蕴含 \(aub<avb\)，则称 \(<\) 为\term[kerongxu-zixu]{可容许字序}[admissible word order]。例如先比较字长，再按固定的字母次序作字典比较，就满足这些条件。
\end{definition}

对一个非零多项式 \(f\)，其最大的字称为\term[shou-danxiangshi]{首单项式}[leading monomial]，记为 \(\operatorname{LM}(f)\)。若其系数为一，称 \(f\) 首一。规定一族规则 \(\ell_\alpha\to r_\alpha\)，其中 \(\ell_\alpha\ne1\)，右侧为有限线性组合且其中每个字都严格小于 \(\ell_\alpha\)。这相当于使用首一关系 \(g_\alpha=\ell_\alpha-r_\alpha\)。

\ProseParagraphBreak
一次重写将多项式中一个完整项 \(c\,a\ell_\alpha b\) 换成 \(c\,ar_\alpha b\)，然后合并同类项。字若不含任何 \(\ell_\alpha\) 作为连续子字，称为\term[zhenggui-zi]{正规字}[normal word]。右侧的字小于 \(\ell_\alpha\)，还必须在加入任意前缀 \(a\)、后缀 \(b\) 后继续小于 \(a\ell_\alpha b\)。可容许字序的双侧相容性正保证这件事；仅有良序而没有这项相容性，局部的下降不能保证大字内部的重写也下降。

\ProseParagraphBreak
同一字中若有两处可替换，可能出现两类\term[chongxie-qiyi]{歧义}[rewriting ambiguity]。一类是包含：\(\ell_\alpha=a\ell_\beta b\)，可先改整个左端，也可先改其中一段；另一类是重叠：\(\ell_\alpha=au\)、\(\ell_\beta=ub\)，其中 \(a,u,b\) 都非空，字 \(aub\) 有两种第一步。相同左端的不同规则也算包含歧义。若两种第一步结果能继续重写为同一个多项式，称这项歧义可解。

\ProseParagraphBreak
例如规则 \(x^2\to1\)、\(yx\to0\) 在 \(yxx\) 上重叠：改末尾 \(xx\) 得到 \(y\)，改开头 \(yx\) 却得到零。规则 \(xyx\to x\)、\(yx\to0\) 则在 \(xyx\) 上包含：改整个字得到 \(x\)，改其中的末尾 \(yx\) 得到零。两组原始规则都下降，但各自的两支已不能继续合到一起。\cref{b2:prop:diamond-ambiguity-failures}将计算这两组关系实际定义的商代数。

\ProseParagraphBreak
两处彼此分离时不形成新的歧义：在字 \(a\ell_\alpha b\ell_\beta c\) 中，先改哪一处，再改另一处，都会得到 \(ar_\alpha br_\beta c\)。这里利用有限和的分配律，不交换任何字母。

\subsection{Diamond 引理与正规字}
\label{b2:subsec:2006:02}

\begin{theorem}[Diamond 引理，首一字重写版本]\label{b2:thm:diamond-lemma}
设 \(k\) 为域，\(X\) 为有限字母集，\(<\) 为 \(X^*\) 上的可容许字序。给定一族首一字重写规则 \(\ell_\alpha\rightarrow r_\alpha\)，其中每个 \(\ell_\alpha\) 都是非空字，\(r_\alpha\) 是严格小于 \(\ell_\alpha\) 的字的有限 \(k\) 线性组合。每次重写在多项式中将一个完整项 \(c\,a\ell_\alpha b\) 换为 \(c\,ar_\alpha b\)，其中 \(c\in k\)，\(a,b\in X^*\)，然后合并同类项，则从任意多项式出发的每个重写过程都在有限步后停止。若全部包含与重叠歧义可解，则每个多项式有唯一的正规字线性组合形式；这里正规字指不含任何 \(\ell_\alpha\) 为连续子字的字。令 \(I\) 为全部 \(\ell_\alpha-r_\alpha\) 生成的双边理想，则正规字的商类构成 \(k\langle X\rangle/I\) 的一组 \(k\) 基。
\end{theorem}

这项结果称为\term[Diamond-yinli]{Diamond 引理}[Diamond lemma]。这里证明的是自由结合代数上采用首一规则与可容许全序的专门形式；原论文还讨论更一般的约化系统\cite{Bergman1978DiamondLemma}。只写若干等式而不给下降次序，还不能使用它。

\begin{proof}
每次重写将一个完整项换为有限多个更小的字，下降依据是良序，右侧的有限性则保证仍得到有限多项式。当前最大字未必每步都下降，因为也可以先改较小项；要证明任意选择都终止，对有限多项式的最大字 \(w\) 作良序归纳。只要 \(w\) 尚未被替换，对较小项进行的操作都是从一个有限的较小字多项式出发，按归纳假设不可能无限进行。较小字也不能产生 \(w\)。若 \(w\) 可替换，把它的整个系数项换掉后，余下所有字都小于 \(w\)，再次由归纳假设终止；若 \(w\) 不可替换，只需处理下面各项。最小字一没有可用规则，提供归纳起点。合并同类项造成的消去只会删除项，不会破坏这个下降论证。

再对单个字 \(w\) 作良序归纳证明正规结果唯一，并记为 \(N(w)\)。假设所有更小的字已经有唯一正规结果，将 \(N\) 线性延伸到它们的有限和上。若 \(w\) 正规，规定 \(N(w)=w\)。否则，选择一次替换得到多项式 \(f\)，其各字小于 \(w\)，候选结果为 \(N(f)\)。须证明它不依赖第一步。

比较两种第一步。如果两处互不相交，继续改另一处会得到上一小节写出的同一个表达式，因此按归纳假设两者的正规结果相同。如果两处包含或重叠，它们在某个连续子字中构成一项题设歧义。该歧义的两支能化为共同结果；在两侧加上原字的相同前缀和后缀，仍是合法的重写序列。所有这些字在第一步以后都严格小于 \(w\)，所以归纳假设再次保证两支的正规结果相同。这就排除了第一步的全部选择，并完成归纳。

由线性延伸，得到到正规字线性生成空间的映射 \(N\)。每次重写都保持 \(N\)，所以每个终止结果等于 \(N(f)\)。又每次重写之差属于 \(I\)，因而 \(f-N(f)\in I\)，这已经保证正规字的商类生成商代数。为证明独立，还要排除正规字之间藏在 \(I\) 中的非零关系。对任意字 \(a,b\) 和规则，有 \(N(a\ell_\alpha b)=N(ar_\alpha b)\)，所以 \(N\) 零化 \(a(\ell_\alpha-r_\alpha)b\)。这些元素线性生成 \(I\)，故 \(N(I)=0\)，即 \(I\subseteq\ker N\)；而 \(N(f)=0\) 时，\(f=f-N(f)\in I\)，给出反向包含。因此 \(\ker N=I\)。正规字上的 \(N\) 是恒等映射，所以它们的非零线性组合不可能属于 \(I\)，得到独立性。
\end{proof}

例如选择 \(y>x\) 的长度字典序，规则 \(yx\to xy+1\) 没有自身的包含或重叠歧义：左端的非空真后缀为 \(x\)，非空真前缀为 \(y\)，它们不同。正规字恰为 \(x^iy^j\)，所以引理给出 Weyl 代数的正规单项式基的另一证明。规则 \(yx\to qxy\) 同样处理量子平面；这里依赖的是字序和歧义检查。

\subsection{PBW 型代数的有效约化}
\label{b2:subsec:2006:03}

当规则有限且系数可精确计算时，可以固定一种约化程序：每步取最大的可约字，再选取其中最左的可约位置；若仍有多条规则，按事先编号选择。上一定理保证停止。若歧义条件也成立，输出就与这个具体选择无关，并且
\[
f\in I\quad\Longleftrightarrow\quad N(f)=0.
\]
因此正规式可以用于相等判定与理想成员判定。终止性本身没有给出统一的时间复杂度界。

\ProseParagraphBreak
例如 \(yx\to xy+x^2\) 在 \(y>x\) 的长度字典序下严格下降，且仍没有自身歧义，因此 \(x^iy^j\) 也构成正规单项式基。计算得
\[
yx^2=x^2y+2x^3,\qquad
y^2x=xy^2+2x^2y+2x^3.
\]
这个代数也可用\cref{b2:thm:ore-extension-construction}构造：取 \(R=k[x]\)、\(\sigma=\operatorname{id}\)、\(\delta(f)=x^2f'\)，再把扩张变量记为 \(y\)。两种方法一个检查重写歧义，一个通过导子建立乘法，得到同一条关系及同一组基。

\begin{definition}\label{b2:def:noncommutative-groebner-basis}
设 \(k\) 为域，\(X\) 为有限字母集，\(<\) 为 \(X^*\) 上的可容许字序，\(I\subseteq k\langle X\rangle\) 为双边理想，\(G\subseteq I\setminus\{0\}\) 为 \(I\) 的一组生成元。若每个非零 \(f\in I\) 的首单项式都含某个 \(g\in G\) 的首单项式为连续子字，则称 \(G\) 为 \(I\) 关于 \(<\) 的\term[feijiaohuan-Groebner-ji]{非交换 Gröbner 基}[noncommutative Gröbner basis]。
\end{definition}

满足 Diamond 引理条件的首一关系族就是这样的基。事实上，若 \(f\in I\) 的最大字不可约，它在任何对较小字的重写中都不可能被消去，所以 \(N(f)\ne0\)，与 \(N(I)=0\) 矛盾。

\ProseParagraphBreak
若某项歧义产生两个不同的正规结果，它们之差是理想中的新关系；可以将差的首项归一化，加入新的下降规则，再检查新出现的歧义。这种补全同样利用尚未消去的关系来扩大首项条件，但非交换字的歧义来自连续包含与连续重叠，不能用交换单项式的最小公倍数替代。即使原规则有限，每次新增的左端仍可能产生下一项歧义，不能由单次约化会停止推断整个补全过程也会停止。

\subsection{非交换 Gröbner 基的有限性问题}
\label{b2:subsec:2006:04}

有限条原始关系不保证只需补充有限多条规则，甚至不保证在给定字序下存在任何有限 Gröbner 基。下面对一个仅有两个生成元、一条关系的代数证明这一点：不仅所写的补全过程无穷延续，其首字反链还将排除所有有限替代。不过这个结论依赖指定的字序，换序后同一理想可以有有限 Gröbner 基。

\ProseParagraphBreak
在 \(k\langle x,y\rangle\) 中，取 \(x>y\) 的长度字典序，从原关系 \(x^2-xy\) 得到规则 \(x^2\to xy\)。对同一个字 \(x^3\)，两种第一步产生分叉：
\[
(x^2)x\longrightarrow xyx,
\qquad
x(x^2)\longrightarrow x^2y\longrightarrow xy^2.
\]
两支的末项都已不能用原规则继续约化，却应代表同一个商类。为使这项歧义可解，须补上 \(xyx\to xy^2\)。

\ProseParagraphBreak
这种补充会继续下去。若已补到规则 \(xy^nx\to xy^{n+1}\)，它与原规则在 \(xy^nx^2\) 上又有两种约化：
\[
\begin{aligned}
(xy^nx)x&\longrightarrow xy^{n+1}x,\\
xy^n(x^2)&\longrightarrow xy^nxy\longrightarrow xy^{n+2}.
\end{aligned}
\]
于是下一条规则应是 \(xy^{n+1}x\to xy^{n+2}\)。这些左端中的两个 \(x\) 相隔越来越多的 \(y\)，以前的规则不能直接改写新的左端。下面证明这一过程产生的全部关系确实属于原理想，并形成无限 Gröbner 基。

\begin{proposition}\label{b2:exm:infinite-noncommutative-groebner}
设 \(k\) 为域，在 \(k\langle x,y\rangle\) 中令 \(I=(x^2-xy)\)。采用 \(x>y\) 的长度字典序时，该理想没有有限非交换 Gröbner 基。一组无限首一 Gröbner 基为
\[
g_n=xy^nx-xy^{n+1},\qquad n\ge0.
\]
商代数的正规字为 \(y^a\) 及 \(y^axy^b\)、\(a,b\ge0\)。改用 \(y>x\) 的长度字典序时，单条关系 \(xy-x^2\) 就是有限首一 Gröbner 基，相应正规字为 \(y^ax^b\)、\(a,b\ge0\)。
\end{proposition}
\begin{proof}
\(g_0=x^2-xy\) 是原关系，而直接展开给出
\[
g_{n+1}=g_ny-g_nx+xy^n g_0.
\]
由归纳，全部 \(g_n\in I\)，所以它们生成的双边理想仍恰为 \(I\)。规则 \(xy^nx\to xy^{n+1}\) 保持字长，最后一个 \(x\) 改成较小的 \(y\)，因此下降。

不同左端不互相包含，因为它们各有且只有两个 \(x\)。非平凡重叠只能共享这个末端与另一个左端的首个 \(x\)，得到字 \(xy^axy^bx\)。先改左边，得到 \(xy^{a+b+1}x\)，再改为 \(xy^{a+b+2}\)；先改右边，得到 \(xy^axy^{b+1}\)，再改左边，也得到 \(xy^{a+b+2}\)。因此所有歧义可解。\cref{b2:thm:diamond-lemma}给出正规字基，而不可约字恰好是至多含一个 \(x\) 的字，即所列形式。

这也确定了理想的全部首单项式：它们恰为包含某个 \(xy^nx\) 的字。一个方向由每个 \(g_n\in I\) 及两边乘字得到；另一个方向由正规式判定得到，因为理想中的非零元素不可能有不可约的最大字。

各个字 \(xy^nx\) 对连续子字包含关系构成无限反链，因而理想的这些最小首字不能由有限多个其他首字统一覆盖。具体地，若存在有限 Gröbner 基，对每个 \(n\)，其中某个首单项式必须是 \(xy^nx\) 的子字；而它又必须包含某个 \(xy^mx\)。因为两个 \(x\) 必须同时保留，只能有 \(m=n\)，且该首单项式就是整个 \(xy^nx\)。于是基中必须为每个 \(n\) 提供一个不同的首单项式，与有限性矛盾。

改用 \(y>x\) 的长度字典序后，\(xy-x^2\) 的首字为 \(xy\)，规则 \(xy\to x^2\) 严格下降。它没有非平凡的自身包含歧义，也没有重叠歧义，因为唯一的非空真后缀 \(y\) 与唯一的非空真前缀 \(x\) 不同。因此\cref{b2:thm:diamond-lemma}给出正规字基，并由前面首项判定说明这条关系构成 Gröbner 基。不含连续子字 \(xy\) 的字只能先写若干个 \(y\)，再写若干个 \(x\)，恰为 \(y^ax^b\)。这证明同一理想的有限 Gröbner 基存在性确实可以随字序改变。
\end{proof}

上述理想在 \(x>y\) 字序下的无限规则仍可用于计算：给定长度为 \(N\) 的字，只可能匹配 \(0\le n\le N-2\) 的左端 \(xy^nx\)。沿字从左到右寻找相邻两个 \(x\) 之间全由 \(y\) 组成的连续子字，就能判定可约位置并执行下一步；对有限多项式逐项检查即可。每步都在可容许字序下下降，所以在系数运算有效时得到终止的正规式计算。“没有有限 Gröbner 基”不等于“任何计算都不可能”。这里证明有限展示不保证给定字序下有有限 Gröbner 基，同时另行证明了这组无限规则可以有效匹配；不能仅由有限展示推得后一性质。第一卷附录F已经给出交换 Gröbner 基的有限性及基本算法；第三卷第8章进一步发展交换情形的计算与几何应用，第三卷附录F则处理局部标准基。使用时须分别核对对象、字序或项序，以及终止性的实际依据。

\begin{proposition}\label{b2:prop:diamond-ambiguity-failures}
设 \(k\) 为域。在 \(k\langle x,y\rangle\) 中，规则 \(x^2\to1\)、\(yx\to0\) 都按长度字典序下降，却在 \(yxx\) 上有不可解的重叠歧义，且
\[
k\langle x,y\rangle/(x^2-1,yx)\cong k[s]/(s^2-1).
\]
在这个商中 \(y=0\)，基为 \(1,x\) 的商类。规则 \(xyx\to x\)、\(yx\to0\) 也都下降，却在 \(xyx\) 上有不可解的包含歧义，且
\[
k\langle x,y\rangle/(xyx-x,yx)\cong k[s].
\]
在这个商中 \(x=0\)，基为 \(1,y,y^2,\ldots\) 的商类。
\end{proposition}
\begin{proof}
第一组规则在 \(yxx\) 上，一条约化先改 \(xx\) 而得到 \(y\)，另一条先改 \(yx\) 而得到零；\(y\) 与零对于原规则都已正规，所以歧义不可解。令 \(A=k\langle x,y\rangle/(x^2-1,yx)\)，暂将生成元的商类仍记为 \(x,y\)。结合律与两条关系给出
\[
y=yx^2=(yx)x=0.
\]
向 \(k[s]/(s^2-1)\) 发送 \(x\mapsto s\)、\(y\mapsto0\)，两条关系均变成零，故得到同态 \(A\to k[s]/(s^2-1)\)。反向以 \(s\mapsto x\) 定义同态，因 \(x^2=1\) 而良定义。两个复合保持生成元；在 \(A\) 中，已经证明 \(y=0\)，所以这两个同态互逆。首一多项式 \(s^2-1\) 的带余除法给出 \(1,s\) 的商类为基，得到 \(A\) 中所列的基。这个论证在特征二时也成立。

第二组规则中，改整个 \(xyx\) 得到 \(x\)，改末尾 \(yx\) 却得到零；\(x\) 与零对原规则仍都正规。令 \(B=k\langle x,y\rangle/(xyx-x,yx)\)，其关系给出
\[
x=xyx=x(yx)=0.
\]
发送 \(x\mapsto0\)、\(y\mapsto s\) 得到同态 \(B\to k[s]\)，反向以 \(s\mapsto y\) 得到同态 \(k[s]\to B\)。两个复合都保持生成元，其中 \(B\) 内的 \(x=0\)，所以它们互逆。\(k[s]\) 的幂基给出所列商类为基。两组失败都产生了原先终止结果未显示的额外关系，因此原规则的终止性没有保证正规字的商类线性独立。
\end{proof}
